\section{Трение}

\subsection{Поверхностное трение}

\textbf{Поверхностное трение} возникает при соприкосновении твердых тел. Она \textit{не зависит от скорости} скольжения ($F(v) = \text{const}$):

\begin{center}
\begin{tblr}{
    colspec = {X[c]X[l]},
    vlines,
    hlines,
    row{1} = {font=\bfseries, mode=text},
    row{odd} = {bg=gray!10},
    cells = {m},
    rowsep = 6pt
}
Формула & Описание \\
$F_{\text{пок}} \le \mu N$ & \textit{Трение покоя:} сила, удерживающая тело на месте. Равна внешней силе, пока та не превысит порог \\
$F_{\text{ск}} = \mu N$ & \textit{Трение скольжения:} сила, действующая при движении тела. Направлена против скорости \\
$F_{\text{кач}} = \dfrac{\delta}{R} N$ & \textit{Трение качения:} возникает при деформации поверхности качения;\newline $\delta$ — коэффициент трения качения, $R$ — радиус колеса \\
\end{tblr}
\end{center}

\textbf{Коэффициент трения ($\mu$)} --- безразмерная величина, зависящая от материала и качества обработки двух поверхностей.

\subsection{Сопротивление в жидкостях и газах}
При движении в жидкой или газообразной среде сила трения покоя отсутствует, а сила сопротивления зависит от скорости.

\subsubsection{Вязкостное трение}
Возникает при \textbf{малых скоростях}, когда слои среды плавно обтекают тело без перемешивания:
\[ \vec{F} \sim \vec{v} \]

\subsubsection{Сопротивление среды}
Возникает при \textbf{больших скоростях}, когда за телом образуются вихри:
\[ \vec{F} \sim v^2 \]